This work is on a filtering pipeline for two-dimensional indoor localization with synthetic LiDAR pose measurements and fixed radio beacons. The estimator is an error-state extended Kalman filter whose nominal state is a five-degree-of-freedom 2D NavState matrix containing heading, position, and velocity. Prediction uses noisy body-frame acceleration and yaw rate, while measurement updates use beacon ranges with either fixed or range/RSSI-adaptive covariance and synthetic LiDAR pose measurements with scan-quality-dependent covariance. Three generated indoor scenarios are evaluated: an open hallway with side rooms, a multi-room map, and a larger dense map. Across the results, dead reckoning gives a mean position RMSE of 7.398 m, fixed-covariance beacon fusion gives 0.326 m, adaptive beacon fusion gives 0.286 m, LiDAR-only fusion gives 0.0968 m, and LiDAR plus adaptive beacons gives 0.0891 m. This shows that the adaptive beacon model improves aggregate beacon-only localization relative to fixed covariance, while the synthetic LiDAR pose update is the dominant correction source in the tested scenarios.
